% Source: https://artofproblemsolving.com/community/c3h2432615_100th_msm_post
% Stable identity: aops-topic-2432615-q05
A number made of $k$ twos is $\dfrac{2\left(10^k-1\right)}{9}$, so the sum is
\[N=\dfrac29\sum_{k=1}^{222}\left(10^k-1\right)=\dfrac29\left(\dfrac{10^{223}-10}{9}-222\right)=\dfrac{2\left(10^{223}-2008\right)}{81}.\]
\textbf{The decimal digits of $\dfrac{2}{81}$.} Since $81\cdot12345679=999999999=10^9-1$, we get $\dfrac{2}{81}=\dfrac{24691358}{10^9-1}=0.\overline{024691358}$. Let $Q=\left\lfloor\dfrac{2\cdot10^{223}}{81}\right\rfloor$; its digits are the first $223$ decimal digits of $\dfrac{2}{81}$, that is $24$ blocks $024691358$ followed by $0246913$ (the leading $0$ is not written).

\textbf{The remainder.} As $10^9\equiv1\pmod{81}$ and $223=9\cdot24+7$, we have $10^{223}\equiv10^7\pmod{81}$, and $10^7=81\cdot123456+64$. So $2\cdot10^{223}\equiv128\equiv47\pmod{81}$, i.e. $2\cdot10^{223}=81Q+47$. Therefore
\[N=\dfrac{81Q+47-4016}{81}=Q-\dfrac{3969}{81}=Q-49.\]
\textbf{Digit sum.} One block $024691358$ has digit sum $38$, so $Q$ has digit sum $24\cdot38+(0+2+4+6+9+1+3)=912+25=937$. The number $Q$ ends in $\dots6913$, and subtracting $49$ changes only these last four digits, to $\dots6864$ (no borrowing beyond them). Their digit sum changes from $6+9+1+3=19$ to $6+8+6+4=24$. So the digit sum of $N$ is
\[937+5=\boxed{942}.\]
